\chapter{Sintonización PID mediante Artificial Bee Colony (ABC)}
\label{ch:abc}

\section{Introducción del capítulo}
Artificial Bee Colony es un algoritmo inspirado en la búsqueda cooperativa de alimento por
parte de una colonia de abejas. Su interés para la sintonización PID proviene de la forma en
que reparte el trabajo entre exploración, explotación y reemplazo de soluciones estancadas,
lo que puede ayudar a mantener diversidad poblacional durante la optimización. Esto resulta
atractivo en problemas donde varias combinaciones de ganancias producen respuestas
aparentemente buenas, pero no todas son igual de robustas.

\section{Objetivo del capítulo}
Este capítulo presenta la lógica del ABC y describe cómo puede adaptarse al ajuste del PID en
una microred aislada. La meta es comprender el papel que cumplen las abejas empleadas,
observadoras y exploradoras dentro del proceso de búsqueda, y evaluar por qué esta división
de funciones puede ser útil en una comparación algorítmica amplia.

\section{Inspiración del algoritmo}
El modelo distingue tres tipos de agentes. Las abejas empleadas exploran y refinan fuentes de
alimento ya conocidas; las observadoras seleccionan probabilísticamente las fuentes más
prometedoras para intensificar la búsqueda; y las exploradoras abandonan soluciones pobres o
estancadas para descubrir nuevas regiones del espacio. Esta organización constituye un
mecanismo natural de balance entre memoria, selección y reinicio.

\section{Modelo matemático del ABC}
En términos de optimización, cada fuente de alimento representa una solución candidata. Una
nueva solución se genera perturbando una fuente existente con respecto a otra solución de la
población, lo que introduce variación dirigida sin perder completamente la referencia previa.
Si la nueva solución mejora el costo, reemplaza a la original; si no mejora durante cierto número
de intentos, la fuente se abandona y se crea una nueva solución aleatoria, preservando así la
diversidad del proceso.

\section{Pseudocódigo}
\begin{algorithm}[H]
\caption{Esquema general del ABC para sintonización PID}
\begin{algorithmic}[1]
\State Inicializar fuentes de alimento en el espacio de ganancias del PID
\State Evaluar cada fuente con la función de costo del problema
\While{no se cumpla el criterio de parada}
    \State Fase de abejas empleadas: generar y evaluar soluciones vecinas
    \State Fase de observadoras: seleccionar fuentes promisorias e intensificar búsqueda
    \State Fase de exploradoras: reemplazar fuentes estancadas por nuevas soluciones
    \State Actualizar la mejor solución global
\EndWhile
\State Devolver la mejor tripleta $\left[K_p,K_i,K_d\right]$
\end{algorithmic}
\end{algorithm}

\section{Parámetros de ajuste}
Los parámetros centrales del ABC son el tamaño de la colonia, el número de ciclos y el límite
de intentos sin mejora antes de abandonar una fuente. Este último parámetro es particularmente
importante porque regula cuándo conviene insistir en una región prometedora y cuándo es mejor
reiniciar la búsqueda. Una configuración demasiado conservadora puede frenar la exploración;
una demasiado agresiva puede impedir la explotación adecuada de soluciones valiosas.

\section{Aplicación a la sintonización PID}
En el problema de interés, cada fuente de alimento codifica una combinación de $K_p$, $K_i$ y
$K_d$. La evaluación se realiza a través de simulaciones del sistema hidro-solar ante cambios
de carga e intermitencia solar, usando ITAE y métricas complementarias como tiempo de
asentamiento y sobreimpulso. Así, la comparación con PSO, GWO y Eagle Strategy se mantiene
en un marco homogéneo, donde lo que cambia es la dinámica de búsqueda y no la definición
del problema.

\section{Cierre del capítulo}
Con ABC se completa el bloque de métodos metaheurísticos seleccionados para esta obra. A
partir de aquí, la discusión se desplaza desde el funcionamiento individual de cada algoritmo
hacia el diseño comparativo que permite enfrentarlos bajo criterios y escenarios comunes.
